\documentclass[11pt]{article}

\usepackage[T1]{fontenc}
\usepackage{lmodern}
\usepackage{amsmath,amssymb,amsthm,mathtools}
\usepackage[margin=1in]{geometry}
\usepackage{microtype}
\usepackage[hidelinks]{hyperref}

\newtheorem{lemma}{Lemma}
\newtheorem{proposition}[lemma]{Proposition}
\newtheorem{corollary}[lemma]{Corollary}
\theoremstyle{remark}
\newtheorem{remark}[lemma]{Remark}

\newcommand{\N}{\mathbb N}
\newcommand{\R}{\mathbb R}

\title{Global Angular Blocks for the Middle-Ratio Shell Proxy}
\author{}
\date{Global analytic angular module}

\begin{document}
\maketitle

\begin{abstract}
The global structural reduction for the spherical-shell proxy leaves an
angular rounding defect $\mathcal E_{\rm ang}$ and a smooth positive payment
$\mathcal H$.  This note controls the defect for every $0<\rho<1$ and $K>0$
with at least one active radial layer.  A single paired radial-block argument,
using only the sharp slope cap $-A'\leq1/2$ and the exact strict angular
rounding, gives
\[
 \mathcal E_{\rm ang}
 <\frac{(1-\rho^2)K^2}{6}-\frac N2,
\]
where $N$ is the number of active shifted radial layers.  On the
upward-closed middle-ratio region
\(
 \rho_c\leq\rho\leq39/50,
 \ K\geq k_{11}(\rho)
\), one has $N\geq1$.  The global middle-ratio theorem is therefore reduced
to the smooth scalar inequality \eqref{eq:scalar-gap}, which is established
in the companion scalar-positivity module.
\end{abstract}

\section{Definitions and the target payment}

For $a>0$ and $z\geq0$, define
\begin{equation}
 G_a(z)=
 \begin{cases}
 \displaystyle\frac1\pi\left(\sqrt{a^2-z^2}
       -z\arccos\frac za\right),&0\leq z<a,\\[4pt]
 0,&z\geq a.
 \end{cases}
 \label{eq:G}
\end{equation}
Fix $0<\rho<1$ and $K>0$, and put
\begin{equation}
 A(z):=G_K(z)-G_{\rho K}(z),
 \qquad
 T:=A(0)=\frac{(1-\rho)K}{\pi}.
 \label{eq:A}
\end{equation}
The function $A$ decreases strictly from $T$ to $0$.  Denote its
decreasing inverse by $R:[0,T]\to[0,K]$, and set
\begin{equation}
 F(t):=R(t)^2,
 \qquad
 q(z):=-A'(z),
 \qquad
 \tau:=A(\rho K)=KG_1(\rho).
 \label{eq:inverse}
\end{equation}
Direct differentiation gives
\begin{equation}
 q(z)=\frac1\pi
 \begin{cases}
 \displaystyle
 \arccos\frac zK-\arccos\frac z{\rho K},&0<z<\rho K,\\[6pt]
 \displaystyle\arccos\frac zK,&\rho K\leq z<K.
 \end{cases}
 \label{eq:q}
\end{equation}
Consequently
\begin{equation}
 0\leq q(z)\leq\frac12,
 \qquad
 q\text{ increases on }(0,\rho K),
 \qquad
 q\text{ decreases on }(\rho K,K).
 \label{eq:q-shape}
\end{equation}

Let
\begin{equation}
 x_n:=n-\frac14,
 \qquad
 M(r):=\max\left\{0,\left\lceil r-\frac12\right\rceil\right\},
 \qquad
 N:=\#\{n\geq1:x_n<T\}.
 \label{eq:grids}
\end{equation}
For $1\leq n\leq N$, write
\begin{equation}
 r_n:=R(x_n),
 \qquad m_n:=M(r_n),
 \label{eq:rn-mn}
\end{equation}
and define the exact angular rounding defect
\begin{equation}
 \mathcal E_{\rm ang}(\rho,K)
 :=\sum_{n=1}^N(m_n^2-r_n^2).
 \label{eq:Eang}
\end{equation}
The strict convention is built into $M$: at $r=m+1/2$ one has
$M(r)=m$, not $m+1$.

The smooth payment obtained from the retained radial Stieltjes remainder
is
\begin{equation}
 \begin{aligned}
 \mathcal H(\rho,K):={}&
 \frac{\rho^2K^2}{4}
 -\frac{\pi\rho K}{4(1-\rho)}\\
 &+\frac{2\pi\rho K}{\arccos\rho}
   \left(\frac{\tau}{4(1+2\tau)}+\frac3{32}\right)\\
 &+\int_{\rho K}^{K}
 2z\,\frac{A(z)(1+A(z))}{(1+2A(z))^2}\,dz.
 \end{aligned}
 \label{eq:H}
\end{equation}
The preceding structural reduction proves the direct lower bound
\begin{equation}
 W(\rho,K)-P(\rho,K)
 \geq\mathcal H(\rho,K)-\mathcal E_{\rm ang}(\rho,K).
\label{eq:structural-implication}
\end{equation}
This note derives a smooth upper bound for
$\mathcal E_{\rm ang}$.

\section{A one-layer paired-block estimate}

The cap $q\leq1/2$ in \eqref{eq:q-shape} forces a useful separation:
one unit of radial action spans at least two units of angular radius.  The
shift $x_n=n-1/4$ leaves a left block of length $3/4$, enough to pay the
worst angular rounding in that layer.

\begin{lemma}[One-layer paired block]
\label{lem:one-layer}
For every active $n$,
\begin{equation}
 r_n
 \leq\frac43\int_{n-1}^{x_n}R(t)\,dt-\frac34,
 \label{eq:rn-block}
\end{equation}
and
\begin{equation}
 m_n^2-r_n^2<r_n+\frac14.
 \label{eq:rounding-pointwise}
\end{equation}
Both statements remain valid, with the displayed strictness, at every
radial--angular common wall.
\end{lemma}

\begin{proof}
Fix $t\in[n-1,x_n]$ and put $z=R(t)\geq r_n$.  Since
$A(r_n)=x_n$, \eqref{eq:q-shape} gives
\[
 x_n-t=A(r_n)-A(z)
 =\int_{r_n}^{z}q(s)\,ds
 \leq\frac{z-r_n}{2}.
\]
Hence
\[
 R(t)\geq r_n+2(x_n-t).
\]
Integration over the interval of length $3/4$ yields
\[
 \int_{n-1}^{x_n}R(t)\,dt
 \geq\frac34r_n+\frac9{16},
\]
which is \eqref{eq:rn-block}.

By strict counting, $m_n<r_n+1/2$ even when $r_n$ is a half-integer:
at $r_n=m+1/2$ the value is $m_n=m$.  Squaring gives
\[
 m_n^2<(r_n+1/2)^2=r_n^2+r_n+1/4,
\]
and proves \eqref{eq:rounding-pointwise}.
\end{proof}

\begin{proposition}[Global angular block bound]
\label{prop:global-E}
For every $0<\rho<1$ and every $K>0$ with $N\geq1$,
\begin{equation}
 \boxed{
 \mathcal E_{\rm ang}(\rho,K)
 <\frac{(1-\rho^2)K^2}{6}-\frac N2.}
 \label{eq:global-E-bound}
\end{equation}
In particular,
\begin{equation}
 \mathcal E_{\rm ang}(\rho,K)
 <\frac{(1-\rho^2)K^2}{6}-\frac12.
 \label{eq:smooth-E-bound}
\end{equation}
\end{proposition}

\begin{proof}
The left blocks $[n-1,x_n]$, $1\leq n\leq N$, are disjoint subsets of
$[0,T]$.  Summing \eqref{eq:rn-block} and then
\eqref{eq:rounding-pointwise} gives
\begin{align*}
 \mathcal E_{\rm ang}
 &<\sum_{n=1}^N\left(r_n+\frac14\right)\\
 &\leq\frac43\sum_{n=1}^N
       \int_{n-1}^{x_n}R(t)\,dt-\frac N2\\
 &\leq\frac43\int_0^T R(t)\,dt-\frac N2.
\end{align*}
The area identity for decreasing inverses gives
\[
 \int_0^T R(t)\,dt=\int_0^K A(z)\,dz.
\]
Scaling $z=as$ in \eqref{eq:G} and using
\[
 \int_0^1\sqrt{1-s^2}\,ds=\frac\pi4,
 \qquad
 \int_0^1s\arccos s\,ds=\frac\pi8,
\]
shows that
\[
 \int_0^aG_a(z)\,dz=\frac{a^2}{8}.
\]
Consequently
\[
 \int_0^T R(t)\,dt=\frac{(1-\rho^2)K^2}{8},
\]
which proves \eqref{eq:global-E-bound}.  Equation
\eqref{eq:smooth-E-bound} follows from $N\geq1$.
\end{proof}

Combining the structural and angular estimates removes the staircase defect
from the live handoff altogether.

\begin{proposition}[Direct analytic proxy gap]
\label{prop:direct-proxy-gap}
For every $0<\rho<1$ and $K>0$ with $N\geq1$,
\begin{equation}
 \boxed{
 W(\rho,K)-P(\rho,K)
 >\mathcal H(\rho,K)
  -\frac{(1-\rho^2)K^2}{6}+\frac N2.}
 \label{eq:direct-proxy-gap}
\end{equation}
This one inequality includes the first and terminal radial pieces and every
strict angular wall; no separate common-jump estimate is a premise.
\end{proposition}

\begin{proof}
Insert \eqref{eq:global-E-bound} into
\eqref{eq:structural-implication}.
\end{proof}

\section{The upward-closed middle-ratio scalar inequality}

Set
\begin{equation}
 \rho_c:=\frac1{1+2\pi},
 \qquad
 k_{11}(\rho):=
 \sqrt{\frac{\pi^2}{(1-\rho)^2}+132},
 \label{eq:compact-constants}
\end{equation}
and consider the upward-closed middle-ratio region
\begin{equation}
 \mathcal R_{\rm c}:=
 \left\{(\rho,K):
 \rho_c\leq\rho\leq\frac{39}{50},\quad
 K\geq k_{11}(\rho)\right\}.
 \label{eq:compact-region}
\end{equation}
On this set,
\begin{equation}
 T\geq\frac{(1-\rho)k_{11}(\rho)}\pi
 =\sqrt{1+\frac{132(1-\rho)^2}{\pi^2}}>1,
 \label{eq:N-positive}
\end{equation}
so $N\geq1$.  Proposition~\ref{prop:direct-proxy-gap} reduces the global
middle-ratio theorem to the single inequality
\begin{equation}
 \boxed{
 \mathcal S(\rho,K):=
 \mathcal H(\rho,K)
 -\frac{(1-\rho^2)K^2}{6}
 +\frac12>0
 }
 \qquad ((\rho,K)\in\mathcal R_{\rm c}).
 \label{eq:scalar-gap}
\end{equation}
Everything in $\mathcal S$ is smooth at every finite point of
$\mathcal R_{\rm c}$; the floor functions and all radial--angular common
jumps have disappeared.  Substituting \eqref{eq:H} gives the explicit form
\begin{equation}
 \begin{aligned}
 \mathcal S(\rho,K)={}&
 \frac{(5\rho^2-2)K^2}{12}
 -\frac{\pi\rho K}{4(1-\rho)}+\frac12\\
 &+\frac{2\pi\rho K}{\arccos\rho}
   \left(\frac{KG_1(\rho)}{4(1+2KG_1(\rho))}
         +\frac3{32}\right)\\
 &+\int_{\rho K}^{K}
 2z\,\frac{A(z)(1+A(z))}{(1+2A(z))^2}\,dz.
 \end{aligned}
 \label{eq:scalar-expanded}
\end{equation}
The elementary identity
\begin{equation}
 \frac{u(1+u)}{(1+2u)^2}
 =\frac14-\frac{1}{4(1+2u)^2}
 \label{eq:kernel-loss}
\end{equation}
turns the same scalar gap into the more useful loss form
\begin{equation}
 \begin{aligned}
 \mathcal S(\rho,K)={}&
 \frac{(1+2\rho^2)K^2}{12}
 -\frac12\int_{\rho K}^{K}
       \frac{z}{(1+2A(z))^2}\,dz\\
 &-\frac{\pi\rho K}{4(1-\rho)}+\frac12
 +\frac{2\pi\rho K}{\arccos\rho}
   \left(\frac{KG_1(\rho)}{4(1+2KG_1(\rho))}
         +\frac3{32}\right).
 \end{aligned}
 \label{eq:scalar-loss-form}
\end{equation}
Thus the only non-endpoint quantity still requiring control is the
positive boundary-layer loss integral in
\eqref{eq:scalar-loss-form}.

\begin{remark}[Status and next step]
The direct proxy gap \eqref{eq:direct-proxy-gap} is proved analytically from
one retained-radial estimate and one one-layer angular estimate.  The first
remaining step within this module is the smooth scalar inequality
\eqref{eq:scalar-gap}.  The companion scalar-positivity module proves it on
all of $\mathcal R_{\rm c}$ by a uniform upper bound for the boundary-layer
loss in \eqref{eq:scalar-loss-form} and strict positivity on the base curve.
No upper cutoff in $K$ occurs in the angular estimate or in the verification
that $N\geq1$.
\end{remark}

\end{document}
